\section{The research question and the distinction that matters}
\label{ir:sec:scope}

\subsection{A precise continuation of the manuscript}
The source is the ProveIt manuscript \emph{Polylogarithms and their
Arithmetic Bridges}, inspected at commit
\begin{center}\code{aeafd6e9b843b589a3bf66850d15bf8ad0e196dd}.\end{center}
Its distribution-jets continuation defines independent prime weights and
proves a conductor-character normal form over a characteristic-zero
splitting field. Its research programme explicitly separates integral
Smith forms, torsion, and modular specializations from those results
\cite{repo,distribution-report}. This article addresses that separation.
The source's positive-characteristic and integral questions are not answered
by tensoring a character calculation with a smaller coefficient ring:
character projectors divide by group orders, and this division is precisely
what is unavailable at some primes.

The construction below uses only sums, subtraction, and multiplication by
weights. Consequently it can be specialized to an arbitrary commutative
ring, including a finite field or a ring with nilpotents. Reflection then
becomes a second, genuinely different operation. Even when a module is
free over $\Z$, quotienting it by $c-1$ or $c+1$ for an involution $c$ can
produce torsion. Rational eigenspace dimensions do not see this torsion.

The main concrete outcomes are summarized in Table~\ref{ir:tab:claims}.
The proofs are ordinary algebraic and analytic proofs. Finite exact
computations replay selected consequences and implementation details;
they are not substitutes for the all-level arguments.

\begin{table}[ht]
\centering\small
\begin{tabular}{p{0.25\textwidth}p{0.64\textwidth}}
\toprule
Result & Exact content and scope\\
\midrule
Integral distribution basis & A basis of actual grid symbols over
$\Z[A_p]$, a selected maximal minor equal to $1$, and arbitrary base change.\\[4pt]
Reflection torsion & Both signs have only $2$-torsion, with an explicit
weight-parity criterion and multiplicity.\\[4pt]
Modular geometry & A Koszul description for every characteristic-two
coefficient algebra; an exact field-valued rank locus.\\[4pt]
One-parameter jets & Complete Smith exponents over $k[[t]]$, not merely
a first-order rank calculation.\\[4pt]
Cyclotomic identities & Polynomial certificates with $3$, $4$, and $5$
raw rows; a level-$15$ derivative identity for every order.\\[4pt]
Audit & A precise Clausen-parity terminology correction; integral versus
field-coefficient qualifications for relation calculations.\\
\bottomrule
\end{tabular}
\caption{What this continuation proves. The classical ordinary
specializations are acknowledged separately.}\label{ir:tab:claims}
\end{table}

\subsection{Classical antecedents and limits of novelty claims}
The ordinary universal distribution is classical. Kubert's 1979 papers
study its generators and its $\Z/2\Z$ cohomology
\cite{kubert-distribution,kubert-sign}. In particular the ordinary
specialization of the torsion count below recovers the known sign
cohomology count; it is not a new resolution of a previously unknown
ordinary case. Ouyang develops the broader universal norm distribution
and an Anderson-type resolution \cite{ouyang}. These are important
antecedents for the present finite weighted complex.

The contribution here is a self-contained, explicit treatment of the
manuscript's particular presentation, including its chosen raw-point
basis, universal unit minor, arbitrary-ring specializations, full weighted
reflection calculation, and contact-order Smith law. These statements
supply material absent from the inspected manuscript. We do not assert
that every special case or general structural principle is absent from
the wider distribution literature. The analytic identities are exact
consequences of classical distribution, supplied with short certificates
and all-order pole-cancelled derivatives; they are not advertised as new
functional equations independent of the classical distribution law.

\subsection{Notation and the formal endpoint}
For a positive integer $q$, let
\[
 X_q=(q^{-1}\Z)/\Z,\qquad P(q)=\{p:p\text{ is prime and }p\mid q\},
 \qquad R_q=\Z[A_p:p\in P(q)].
\]
We write $\den(x)$ for the exact order of $x$ in $\Q/\Z$, and set
$\den(0)=1$. A symbol $e_x$ is not itself a special value. It is a generator
in a free module. Under polylogarithm evaluation, $e_0$ will become
$\Li_s(1)=\zeta(s)$. Under Hurwitz evaluation, it becomes
$\zeta(s,1)$, not the undefined expression $\zeta(s,0)$.

For $d\mid q$, put $A_d=\prod_{p\mid d}A_p^{v_p(d)}$.
Weights on composite divisors are thus multiplicative; an unrelated
variable for every divisor would define a different module. The universal
module and its relation submodule are
\begin{align}
 V_q&=\bigoplus_{x\in X_q}R_qe_x,\\
 r_{p,x}&=\sum_{\substack{y\in X_q\\py=x}}e_y-A_pe_x
       \qquad(p\mid q,\ x\in X_{q/p}),\label{ir:eq:prime-row}\\
 \R_q&=\langle r_{p,x}\rangle_{R_q},\qquad
 \D_q=V_q/\R_q.\label{ir:eq:module}
\end{align}
Every sum in \eqref{ir:eq:prime-row} has exactly $p$ terms. For a
commutative ring $S$ and weights $a_p\in S$, write $\D_q(S,a)$ for the
same presentation after $A_p\mapsto a_p$.

\begin{lemma}[Prime rows suffice]\label{ir:lem:divisor}
The relation $r_{d,x}=\sum_{dy=x}e_y-A_de_x$, for $d\mid q$ and
$x\in X_{q/d}$, follows from the prime rows without division.
\end{lemma}
\begin{proof}
For $ab\mid q$ and $x\in X_{q/(ab)}$, direct expansion gives
\[
 r_{ab,x}=\sum_{ay=x}r_{b,y}+A_b r_{a,x}.
\]
Here all $y$ belong to $X_{q/b}$, so every displayed prime or divisor
row is legitimate. The intermediate symbols cancel. Induction on the
number of prime factors of $d$, counted with multiplicity, proves the
claim.
\end{proof}

This elementary observation already prevents a common ambiguity:
``distribution rank'' is meaningful only after the row vocabulary,
endpoint convention, coefficient ring, and divisor-weight convention
have all been fixed.
